\documentclass[conference]{IEEEtran}
\IEEEoverridecommandlockouts
% The preceding line is only needed to identify funding in the first footnote. If that is unneeded, please comment it out.
\usepackage{cite}
\usepackage{booktabs}
\usepackage{amsmath,amssymb,amsfonts}
\usepackage{algorithmic}
\usepackage{graphicx}
\usepackage{textcomp}
\usepackage{xcolor}
\def\BibTeX{{\rm B\kern-.05em{\sc i\kern-.025em b}\kern-.08em
    T\kern-.1667em\lower.7ex\hbox{E}\kern-.125emX}}
\begin{document}

\title{LLM-GridEval: Adaptive Co-Simulation for Smart Grid Security}

\author{\IEEEauthorblockN{ Chenglong Fu}
\IEEEauthorblockA{\textit{\small{Software and Information Systems Dept.}} \\
\textit{\small{Univ. of North Carolina at Charlotte}}\\
\small{Charlotte, USA} \\
chenglong.fu@charlotte.edu}
\and
\IEEEauthorblockN{Shirin Besati}
\IEEEauthorblockA{\textit{\small{Electrical and Computer Engineering Dept.}} \\
\textit{\small{Univ. of North Carolina at Charlotte}}\\
\small{Charlotte, USA} \\
sbesati@charlotte.edu}
\and
\IEEEauthorblockN{ Miao Wang}
\IEEEauthorblockA{\textit{\small{Electrical and Computer Engineering Dept.}} \\
\textit{\small{Univ. of North Carolina at Charlotte}}\\
\small{Charlotte, USA} \\
mwang25@charlotte.edu}
% \and
% \IEEEauthorblockN{4\textsuperscript{th} Given Name Surname}
% \IEEEauthorblockA{\textit{dept. name of organization (of Aff.)} \\
% \textit{name of organization (of Aff.)}\\
% City, Country \\
% email address or ORCID}
% \and
% \IEEEauthorblockN{5\textsuperscript{th} Given Name Surname}
% \IEEEauthorblockA{\textit{dept. name of organization (of Aff.)} \\
% \textit{name of organization (of Aff.)}\\
% City, Country \\
% email address or ORCID}
% \and
% \IEEEauthorblockN{6\textsuperscript{th} Given Name Surname}
% \IEEEauthorblockA{\textit{dept. name of organization (of Aff.)} \\
% \textit{name of organization (of Aff.)}\\
% City, Country \\
% email address or ORCID}
}

\maketitle

\begin{abstract}
Rapid ``API-fication" of electric vehicle (EV) fleets and distributed energy resources (DERs) is transforming the power grid into a software exposed, nonlinearly coupled cyber physical system. Cloud hosted management platforms and programmatic control interfaces create logic level attack surfaces where adversaries can coordinate load altering attacks on EV/DER fleets and induce cascading transmission and distribution (T\&D) effects. Yet most security evaluations for intrusion detection and control based defenses still rely on fixed datasets or pre scripted attack scenarios, leading to an \emph{evaluation validity gap}: a discrepancy between a defense's reported performance against static benchmarks and its true robustness against adaptive, feedback driven adversaries. This paper introduces LLM-GridEval, a formal evaluation framework that embeds an adaptive, tool using large language model (LLM) attacker into a HELICS based T\&D co simulation, exposes a physically constrained attack primitive interface, and enables direct comparison between static and adaptive attackers. The framework models the grid as a multi domain dynamical system, defines attacker policies that act through a schema constrained catalog of primitives, and couples these actions to a co simulation of an IEEE 9 bus transmission system and IEEE 123 node distribution feeders. We instantiate LLM-GridEval in a proof of concept EV capacity manipulation scenario where adaptive attackers issue capacity changes via an EV management and control platform. In this setting, a timing-only adaptive agent underperforms a random baseline, while a strategy-aware adaptive agent matches the baseline's violation duration with fewer actions by diversifying targets and accumulating load. This highlights how attacker adaptivity, modeling assumptions, and domain knowledge can materially affect measured risk, demonstrating that static evaluations may misestimate both the severity and efficiency of adversarial threats in API driven, EV/DER rich grids.
\end{abstract}


\begin{IEEEkeywords}
smart grid security, cyber-physical systems, LLM-based attacks, co-simulation, electric vehicles, intrusion detection, load-altering attacks, HELICS
\end{IEEEkeywords}

\section{Introduction}
\label{sec:introduction}

The digitalization and electrification of transportation are driving a rapid convergence between transmission networks, distribution feeders, and large fleets of software managed electric vehicles (EVs) and distributed energy resources (DERs). Modern grid operations increasingly depend on cloud based management platforms, supervisory control and data acquisition (SCADA), and programmatic interfaces that schedule EV charging and DER setpoints across thousands of devices. While this integration improves observability and flexibility, it also exposes new logic level attack surfaces: adversaries who compromise EV or DER management platforms can coordinate load altering attacks that propagate from distribution feeders into the bulk transmission system, potentially causing line overloads, voltage instability, and large scale load shedding \cite{morris2013icssecurity,cintuglu2017testbeds,yohanandhan2022review,laaSurvey2024}.

To assess grid resilience, the community has developed a rich ecosystem of testbeds, co simulation frameworks, and intrusion detection systems (IDS) for industrial control systems and smart grids \cite{cintuglu2017testbeds,hardy2024helics,yohanandhan2022review}. However, most evaluation methodologies still rely on static datasets or a small number of pre scripted attack scenarios. Widely used traces such as SWaT, WADI, EPIC, and HAI capture fixed attack trajectories under specific operating conditions \cite{swatdataset,wadidataset,epicdataset,haidataset}; many smart grid studies likewise use hand crafted attack scripts that do not change during an experiment \cite{le2020gridattacksim}. We refer to the resulting discrepancy between a defense's reported performance against such static benchmarks and its robustness against reactive adversaries as the \emph{evaluation validity gap}. Static traces and scripts cannot capture the action reaction dynamics of adaptive attackers and can therefore produce misleading confidence in deployed defenses \cite{sahani2023smartgridids}.

Concurrently, advances in large language models (LLMs) and tool using agents have transformed offensive capabilities in IT and industrial control domains. Systems such as PentestGPT, AutoAttacker, and AttackLLM demonstrate that LLM based agents can decompose complex tasks, read documentation, orchestrate external tools, and refine strategies based on feedback \cite{deng2024pentestgpt,autoattacker,attackllm,fang2024llmagents}. In this work, we view LLMs as scalable \emph{cognitive proxies} for human red teams: they provide semantic reasoning and decision making over heterogeneous artifacts (protocols, APIs, logs) and can be embedded into automated evaluation pipelines.

Applying LLM agents to cyber physical systems introduces a key challenge: traditional LLM based red teaming frameworks are largely IT focused and are not aware of physical constraints. In power systems, hallucinated commands---such as impossible voltages, infeasible power injections, or violating device ratings---are not merely incorrect but physically meaningless. To obtain meaningful evaluations, any LLM based attacker for smart grids must therefore be grounded in a \emph{physically constrained primitive interface} that exposes only valid, well typed actions over the grid state, such as EV capacity adjustments within configured bounds.

This paper introduces LLM-GridEval, an evaluation framework that embeds an adaptive, tool using LLM attacker inside a HELICS based transmission and distribution co simulation \cite{hardy2024helics}. LLM-GridEval connects the attacker to the smart grid through a schema constrained catalog of attack primitives (e.g., \texttt{get\_grid\_status()}, \texttt{set\_ev\_capacity}) that are enforced by a sandboxed interface layer. The framework is explicitly designed to measure and analyze the evaluation validity gap by comparing defenses under static, scripted attackers and under adaptive, observation driven policies.

The contributions of this paper are threefold:
\begin{itemize}
	\item \textbf{Theoretical:} We formalize defense evaluation against adaptive attackers in cyber physical power systems by modeling the attacker grid interaction as a partially observable process and by defining the evaluation validity gap between static and adaptive attack processes.
	\item \textbf{Methodological:} We design LLM-GridEval, a framework that combines HELICS based T\&D co simulation with a schema constrained attack primitive catalog and a sandboxed interface layer suitable for LLM based and algorithmic agents.
	\item \textbf{Empirical:} We instantiate an EV capacity load altering case study in which attacker issued EV setpoint changes, applied through the constrained primitive interface in a realistic T\&D co simulation, drive a distribution feeder into overload and solver instability, illustrating how adaptive attackers can expose vulnerabilities that static baselines fail to reveal.
\end{itemize}


\section{Background and Related Work}
\label{sec:background}

\subsection{Smart Grid Co-simulation and Security Datasets}

Smart grid security research relies on cyber-physical testbeds and co-simulation frameworks to study interactions between power system dynamics and control mechanisms \cite{cintuglu2017testbeds,yohanandhan2022review}. HELICS has emerged as a widely used open-source framework for multi-domain co-simulation, supporting time-coordinated interaction between transmission solvers, distribution simulators, and communication models \cite{hardy2024helics}. Security-focused platforms such as GridAttackSim \cite{le2020gridattacksim} combine power system and network simulation but rely on predefined attack scenarios. These platforms represent \emph{Generation~1} evaluation tools: high-fidelity environments with adversary logic fixed in advance.

Similarly, widely used ICS security datasets including SWaT, WADI, EPIC, and HAI \cite{swatdataset,wadidataset,epicdataset,haidataset} capture static recordings of scripted attack campaigns. Survey work highlights that evaluation results from such traces may not generalize to adaptive adversaries \cite{sahani2023smartgridids}. Moreover, most testbeds do not cover EV or DER-rich feeders or load-altering attacks that exploit controllable demand. This motivates evaluation frameworks where dynamic attackers operate \emph{within} co-simulations rather than being pre-scripted around them.

\subsection{Automated Attack Generation and LLM-based Red Teaming}

The security literature has developed models for false data injection attacks (FDIAs) \cite{aoufi2020fdiasurvey} and load-altering attacks \cite{laaSurvey2024}, but these typically assume offline attack planning rather than online adaptation. Automated approaches such as attack graph generation \cite{a2g2v}, control-theoretic red teaming \cite{tackett2024redteam}, and deep reinforcement learning for targeted attacks \cite{targetedAttackSynthesis} represent \emph{Generation~2--3} adversaries that can adapt over time but require scenario-specific state representations.

Recent work applies LLMs to offensive security. PentestGPT \cite{deng2024pentestgpt} and AutoAttacker \cite{autoattacker} demonstrate LLM-based penetration testing in IT systems, while AttackLLM \cite{attackllm} generates attack patterns for ICS testbeds. These \emph{Generation~4} approaches use LLMs as cognitive orchestrators that reason over documentation and call external tools. However, no existing framework (i) couples LLM agents directly into HELICS-style T\&D co-simulation, (ii) exposes physically constrained attack primitives, and (iii) explicitly studies the evaluation validity gap between static and adaptive attackers. LLM-GridEval addresses this gap.

\section{Problem Statement and System/Threat Model}
\label{sec:problem}
%==================================================
\subsection{Problem Statement}

% \begin{figure}
%     \centering
%     \includegraphics[width=1\linewidth]{figures/flow.png}
%     \caption{Conceptual view of defense evaluation with static and adaptive attackers in LLM-GridEval.}
%     \label{fig:eval-flow}
% \end{figure}

%We consider a cyber physical power system instrumented with a defense mechanism \(M\), which can represent an intrusion detection system, an anomaly detector, a controller with protective logic, or a combination of these components. The goal is to evaluate the performance of \(M\) when the system is subjected to adversarial activity. Let \(E(M, A)\) denote an evaluation functional that maps a defense \(M\) and an attack process \(A\) to performance metrics, which can include cyber metrics such as detection rates and latency, physical metrics such as voltage violations and line overload counts, and stability or resilience metrics.

%We model the interaction between the attacker and the defended grid as a partially observable process. At each discrete time step \(t\), the global system state is \(S_t\) and the attacker chooses an action \(A^A_t\) from an attack primitive catalog. Informally, we write this process as a tuple \(\langle S, A, T, R, \Omega, O \rangle\), where \(S\) is the joint cyber physical state space, \(A\) is the attacker action space defined by the primitive catalog, \(T\) is a black box transition implemented by the co simulation \(T(S_t, A^A_t, A^D_t) = \text{Sim}(S_t, A^D_t, A^A_t)\), \(\Omega\) is the space of attacker observations, and \(O: S \rightarrow \Omega\) maps internal system states to observations (e.g., MCP status) available to the attacker. The reward functional \(R\) captures attacker objectives such as overload duration, constraint violations, or defender utility degradation, but is left abstract in this work.

%An attacker policy \(\pi\) selects actions based on observations and knowledge. We write \(A^A_t = \pi(O_t, K, G)\), where \(O_t = O(S_t)\), \(K\) is a knowledge base including protocol specifications and system documentation, and \(G\) is a goal specification encoding objectives such as maximizing overload duration or avoiding detection. Static attackers can be represented as fixed policies \(\pi_{\text{static}}\) in which \(A^A_t\) depends only on time or a predefined script, while adaptive attackers implement policies \(\pi_{\text{adapt}}\) that react to observations and feedback from the co simulation. When \(\pi_{\text{adapt}}\) is implemented by a large language model equipped with tools, the policy can incorporate pattern recognition and symbolic reasoning, as demonstrated in recent LLM based attack agents \cite{deng2024pentestgpt,attackllm,fang2024llmagents}.

%Within this formalism, we define the evaluation validity gap for a defense \(M\) as
%\begin{equation}
%  \Delta E(M) = E(M, A_{\text{static}}) - E(M, A_{\text{adapt}}),
%\end{equation}
%where \(A_{\text{static}}\) denotes the attack process induced by a static policy \(\pi_{\text{static}}\) and \(A_{\text{adapt}}\) denotes the process induced by an adaptive policy \(\pi_{\text{adapt}}\) under comparable resource and access constraints. If \(E\) measures a utility in which larger values indicate better performance, a positive \(\Delta E(M)\) means that the defense appears more effective under static evaluation than under adaptive evaluation, quantifying the evaluation validity gap between traditional benchmarks and realistic, feedback driven attackers.

We consider a cyber physical power system instrumented with a defense mechanism \(M\), which can represent an intrusion detection system, an anomaly detector, and a controller with protective logic. The goal is to evaluate the performance of \(M\) when the system is subjected to adversarial activity. Let \(E(M, A)\) denote an evaluation functional that maps a defense \(M\) and an attack process \(A\) to performance metrics, which can include cyber metrics such as detection rates and latency, physical metrics such as voltage violations and line overload counts, and stability or resilience metrics.
%
An attacker policy \(\pi\) selects actions based on observations and knowledge. Let \(O_t\) denote the attacker-visible observation at time \(t\) (e.g., the subset of system state returned by \texttt{get\_grid\_status()}). We define \(A^A_t = \pi(O_t, K, G)\), where \(K\) is a knowledge base including protocol specifications and system documentation, and \(G\) is a goal specification encoding objectives such as maximizing overload duration or avoiding detection. Static attackers can be represented as fixed policies \(\pi_{\text{static}}\) in which \(A^A_t\) depends only on time or a predefined script, while adaptive attackers implement policies \(\pi_{\text{adapt}}\) that react to observations and feedback from the co-simulation. 

In this work, \(\pi_{\text{adapt}}\) is implemented by an LLM equipped with tools, and the policy can incorporate pattern recognition and symbolic reasoning, as discussed in \cite{deng2024pentestgpt,attackllm,fang2024llmagents}.
%
We define the evaluation validity gap for a defense \(M\) as
\begin{equation}
  \Delta E(M) = E(M, A_{\text{adapt}}) - E(M, A_{\text{static}}),
\end{equation}
where \(A_{\text{static}}\) denotes the attack process induced by a static policy \(\pi_{\text{static}}\) and \(A_{\text{adapt}}\) denotes the process induced by an adaptive policy \(\pi_{\text{adapt}}\) under comparable resource and access constraints. We interpret \(E(M, A)\) as a stress (harm) functional where larger values indicate worse grid performance, such as longer threshold violation duration or higher overload counts. Under this convention, \(\Delta E(M) > 0\) indicates the adaptive attacker induces greater harm than the static baseline, implying that static evaluations underestimate adversarial risk. For example, if \(E\) measures TVD (threshold violation duration in seconds), then \(\Delta E(M) > 0\) reveals that the adaptive attacker achieves longer violation durations than the random baseline, exposing a gap between static and adaptive threat models.



%==================================================
\subsection{System Model and Threat Model}
We model the smart grid as a discrete-time cyber-physical system. The power system combines an IEEE 9-bus transmission network (GridPACK \cite{gridpack}) with IEEE 123-node distribution feeders (GridLAB-D \cite{gridlabd}), integrated via HELICS co-simulation \cite{hardy2024helics}. Six controllable EV charging stations on Feeder A serve as attack targets; two are co-located with storage for islanding capability.

At time step $t$, the system state $S_t = (P_t, C_t)$ comprises physical state (voltages, power injections, switch statuses) and cyber state (SCADA measurements, control application states). The defense mechanism $M$ outputs actions $A^D_t = M(S_{t-k}, \ldots, S_t)$, while the attacker issues actions $A^A_t$ from a primitive catalog. System evolution follows $S_{t+1} = \text{Sim}(S_t, A^D_t, A^A_t)$ via HELICS federation.

The cyber infrastructure includes an EV management and control platform (MCP) server that exposes HTTP endpoints. The threat model assumes an adversary who has compromised the MCP interface, for example through vulnerabilities in backend applications or cloud misconfigurations. The adversary has read access to aggregate feeder measurements and EV setpoints via \texttt{get\_grid\_status()}, and write access to EV charging capacities within configured bounds via \texttt{set\_ev\_capacity(ev\_id, P)}. The adversary cannot directly observe transmission system internal states or protection relay settings. To enforce physical realizability, attack actions are subject to per-device capacity bounds, rate limits, and action budgets. The defender is a threshold-based feeder controller that monitors total power and adjusts EV charging when thresholds are exceeded.


\section{Design of LLM-GridEval}
\label{sec:design}

\subsection{Framework Architecture}
The LLM-GridEval framework integrates the system model into a unified evaluation environment through three layers (Fig.~\ref{fig:framework-architecture}).

\begin{figure*}
    \centering
    \includegraphics[width=0.5\linewidth]{figures/system.png}
    \caption{High level architecture of the LLM-GridEval framework, showing the co-simulation, sandboxed interface, and LLM-based cognitive layers.}
    \label{fig:framework-architecture}
\end{figure*}

\textbf{Co-simulation Layer}: Implements $\text{Sim}(\cdot)$ through HELICS federates with synchronized time advancement: a GridPACK transmission federate (IEEE 9-bus), two GridLAB-D distribution federates (IEEE 123-node feeders; Feeder A contains EV loads, Feeder B contains only background loads), a feeder controller federate implementing threshold-based defense, and an MCP federate connecting the EV management interface.

\textbf{Interface and Tool Layer}: Acts as a sandboxed API boundary between co-simulation and attacker. Each attack primitive is defined by a schema for inputs/outputs, preconditions, and mapping to HELICS operations. The layer performs argument validation and constraint enforcement (e.g., EV capacity bounds, rate limits) before invoking HELICS calls, ensuring all actions are physically realizable and auditable.

\textbf{Cognitive Layer}: Hosts the attacker policy $\pi$ using a ReAct-like interaction loop: (1) observe grid state via \texttt{get\_grid\_status()}, (2) construct prompt with state and available primitives, (3) query LLM for action specification, (4) validate against schema and constraints, (5) execute via sandboxed interface, (6) advance simulation. This closed loop connects LLM reasoning to physically meaningful grid actions while the constrained catalog mitigates hallucinated commands.

This layered separation serves two purposes. First, it enables controlled experimentation: researchers can swap attacker implementations (scripted, RL-based, or LLM-based) while keeping the co-simulation and tool interface fixed, isolating the impact of policy changes. Second, it provides security guarantees: the sandboxed interface ensures that even a fully autonomous LLM agent cannot issue physically invalid commands or bypass rate limits, making automated red-teaming safe for deployment against production-representative grid models.

\subsection{Prototype Implementation}
The prototype instantiates LLM-GridEval for EV capacity attacks on a transmission-distribution co-simulation.

\textbf{Physical Plant}: The IEEE 9-bus transmission system connects to two IEEE 123-node distribution feeders via HELICS value federates with conservative time advancement \cite{hardy2024helics}. Feeder A hosts six EV charging stations (EV1--EV6) distributed across phases A, B, and C at different nodes \cite{ieee123feeder}. EV1 and EV4 are co-located with local storage and can be islanded from the main feeder. Feeder B contains only uncontrollable background loads for demand diversity.

\textbf{MCP Server}: Implemented as a HELICS federate, the MCP server exposes the two attack primitives. \texttt{GET /status} implements \texttt{get\_grid\_status()} and returns feeder power at the swing node, individual EV capacities and actual powers, and islanding states. \texttt{POST /set\_capacity} implements \texttt{set\_ev\_capacity(ev\_id, P)}, validates requests against configured capacity bounds, and publishes setpoints to GridLAB-D. This design mirrors realistic EV fleet management systems where cloud-hosted services provide programmatic control interfaces.

\textbf{Feeder Controller}: The defender is a rule-based controller that samples total feeder power $P_{\text{feeder}}(t)$ at the feeder-head (swing node). When $P_{\text{feeder}} \geq P_{\max}$ (4.2~MW), the controller sheds EV load; when $P_{\text{feeder}} \leq P_{\min}$ (2.6~MW), it restores EV capacities toward nominal levels. This threshold-based control provides a baseline defense for evaluation.

\textbf{LLM Orchestration}: At configurable simulation intervals, a coordination module queries the MCP server via \texttt{get\_grid\_status()} and formats the response into a structured prompt. The prompt includes the current grid state, available primitives with their schemas, attack constraints, and high-level objectives. The LLM returns a JSON action specification that is validated against the primitive schema before execution. For the experiments in this paper, we use a local LLM endpoint (Llama-3.1-8B) to ensure reproducibility and avoid API rate limits during extended campaigns.



\section{Proof-of-Concept Evaluation}

We instantiate LLM-GridEval in an EV load-altering attack scenario to demonstrate the framework's ability to compare static and adaptive attackers under identical conditions. The evaluation addresses two questions: (1) Can adaptive attackers exploit timing patterns to outperform random baselines? (2) What domain knowledge must LLM-based attackers possess to be effective?

\subsection{Experimental Setup}

\subsubsection{Grid Configuration}
The testbed uses the IEEE 123-node distribution feeder connected to an IEEE 9-bus transmission system via HELICS co-simulation. Six EV charging stations (EV1 through EV6) are distributed across phases A, B, and C at different nodes. The feeder operates with a base load of approximately 2.8 to 3.2 MW during the test period (simulation hour 7), with a protection threshold of $P_{\max} = 4.2$ MW. When exceeded, the defensive controller sheds EV load to restore safe operating conditions.

\subsubsection{Defensive Controller}
The threshold-based feeder controller polls total feeder power every 10 seconds. When $P_{\text{feeder}} \geq P_{\max}$, the controller commands all EV stations to zero power. When load returns below 2.6 MW, the controller gradually restores EV charging. This simple defense represents a baseline congestion management mechanism.

\subsubsection{Attack Constraints}
To create meaningful strategic trade-offs, we impose a scarce attack budget on all attacker variants:
\begin{itemize}
    \item \textbf{Cooldown}: Minimum 90 seconds between consecutive attacks (upper-bounding the rate to $\leq$40 attacks/hour)
    \item \textbf{Power bounds}: Each attack sets EV power within [500, 1500] kW
    \item \textbf{Ramping}: Power changes are rate-limited to 100 kW/s to prevent solver instability
\end{itemize}

These constraints ensure that each attack decision carries strategic weight: poor timing or target selection wastes scarce attack opportunities.

The 5-minute experiment window was selected to capture multiple attack-response cycles while keeping computational costs tractable for parameter sweeps. Each experiment begins at simulation hour 7, corresponding to moderate evening load conditions where the base load (2.8--3.2~MW) leaves headroom below the 4.2~MW threshold but is close enough that coordinated EV manipulation can trigger violations.

\subsection{Attack Variants}

We evaluate three attacker policies through the same MCP primitive interface, ensuring that differences in outcomes reflect policy quality rather than interface advantages.

\subsubsection{Random Baseline ($\pi_{\text{static}}$)}
The random attacker represents a static, non-adaptive policy. At each observation interval (5 seconds), if the cooldown has elapsed, the attacker attacks with 30\% probability. Target EV and power level are selected uniformly at random from the valid ranges. This baseline has no timing intelligence and no strategic planning.

\subsubsection{AI-V1: Timing-Only ($\pi_{\text{adapt}}^{(1)}$)}
The first adaptive policy uses an LLM with two-level timing intelligence:

\textbf{Macro-timing} assesses grid stress by computing the headroom between current load and threshold. A macro score of 100 indicates the threshold is already exceeded; scores above 70 indicate favorable attack conditions with headroom below 1 MW.

\textbf{Micro-timing} tracks the defender's observation cycle. The controller acts every 10 seconds; an attack issued immediately after a controller action (cycle position $\approx 0.0$) has nearly 10 seconds before the next defensive response, while an attack at cycle position 0.9 has only 1 second. The micro score maps cycle position to a 0 to 100 scale favoring early-cycle attacks.

The AI-V1 policy waits until the micro-timing score exceeds 70 before attacking, but receives no guidance on target selection or power accumulation dynamics.

\subsubsection{AI-V2: Timing with Strategy ($\pi_{\text{adapt}}^{(2)}$)}
The second adaptive policy extends V1 with explicit strategic knowledge. The LLM prompt includes:
\begin{itemize}
    \item \textbf{Power accumulation model}: Explanation that EV contributions are additive across different targets but overwrite when the same EV is attacked repeatedly
    \item \textbf{Diversification priority}: Instruction to attack unattacked EVs before repeating any target
    \item \textbf{Attack history}: Runtime context showing which EVs have been attacked and their current power levels
\end{itemize}

This additional context enables the LLM to reason about cumulative grid impact rather than optimizing individual attack timing in isolation.

\subsection{Power Accumulation Model}

A critical insight that emerged during V1 testing is that EV power contributions follow an additive-but-overwriting model:
\begin{equation}
P_{\text{total}} = P_{\text{base}} + \sum_{i=1}^{6} P_{\text{EV}_i}
\end{equation}

When an attacker issues a command to EV$_i$, it sets $P_{\text{EV}_i}$ to the specified value, overwriting any previous setpoint for that EV. Consequently, three attacks on EV1 at 1000, 1200, and 1500 kW contribute only 1500 kW total, while three attacks on EV1, EV2, and EV3 at 1500 kW each contribute 4500 kW. This property makes target diversification essential for effective load accumulation. We report the final accumulated EV setpoint $\sum_i P_{\text{EV}_i}$ at the end of each run; repeated attacks on the same EV overwrite its setpoint rather than add.

\subsection{Results}

Table~\ref{tab:results} summarizes results from 5-minute experiments under identical grid conditions and constraints. Each attacker variant was executed three times with different random seeds; we report mean values. The primary metric is Threshold Violation Duration (TVD): cumulative seconds during which feeder load exceeds $P_{\max}$. Higher TVD indicates more effective attacks from the adversary's perspective, and correspondingly, greater stress on the defensive controller. Attack success rate is the fraction of executed attacks after which $P_{\text{feeder}}$ exceeds $P_{\max}$ within the subsequent controller observation interval.

\begin{table}[t]
\centering
\caption{Comparison of Attack Strategies}
\label{tab:results}
\begin{tabular}{lccc}
\toprule
\textbf{Metric} & \textbf{Random} & \textbf{AI-V1} & \textbf{AI-V2} \\
\midrule
\multicolumn{4}{l}{\textit{Primary Metrics}} \\
TVD (seconds) & 240 & 120 & 240 \\
Total attacks executed & 4 & 3 & 3 \\
Attack success rate & 25\% & 33\% & 33\% \\
\midrule
\multicolumn{4}{l}{\textit{Timing Metrics}} \\
Avg. micro-timing score & 75 & 91 & 100 \\
Avg. cycle position & 0.25 & 0.09 & 0.00 \\
\midrule
\multicolumn{4}{l}{\textit{Strategic Metrics}} \\
Unique EVs targeted & 4 & 1 & 3 \\
Avg. power per attack (kW) & 865 & 1133 & 1500 \\
Final accumulated EV setpoint (kW) & 3459 & 1500 & 4500 \\
\bottomrule
\end{tabular}
\end{table}

\subsubsection{AI-V1 Analysis}
Despite achieving superior timing metrics (micro-score 91 vs. 75, cycle position 0.09 vs. 0.25), AI-V1 produced the worst TVD outcome at 120 seconds, half that of the random baseline. Examination of the attack log revealed that V1 attacked EV1 exclusively across all three attacks, with power levels of approximately 700, 1200, and 1500 kW. Each subsequent attack overwrote the previous setpoint, resulting in only 1500 kW of cumulative contribution despite expending three attacks.

This result demonstrates that timing optimization alone is insufficient. The LLM optimized a local objective (attack at good timing) without understanding the global objective (maximize cumulative load). From the LLM's perspective, EV1 was always a valid target, and attacking it at optimal timing appeared to be a good decision.

\subsubsection{AI-V2 Analysis}
With explicit strategic guidance, AI-V2 achieved both perfect timing (micro-score 100, cycle position 0.00) and effective diversification. The attack sequence targeted EV1, then EV2, then EV3, each at maximum power (1500 kW). This produced 4500 kW of cumulative EV load, significantly exceeding the random baseline's 3459 kW.

AI-V2 matched random's TVD of 240 seconds while using 25\% fewer attacks (3 vs. 4). The per-attack efficiency of 80 seconds TVD per attack for V2 versus 60 seconds for random reflects V2's superior resource utilization. Under the scarce budget constraints, this efficiency advantage would compound over longer attack campaigns.

\subsubsection{Ceiling Effect}
Both random and AI-V2 achieved identical TVD despite V2's superior metrics. This ceiling effect arises because the grid at Hour 7 is sufficiently stressed that even suboptimal attacks can trigger violations. Under harder scenarios, such as higher thresholds, lower base load, or faster controller response, the evaluation validity gap $\Delta E(M)$ between adaptive and static attackers would widen.

\subsection{Discussion}
These experiments yield three insights for LLM-based cyber-physical attackers. First, tactical excellence does not guarantee strategic success: AI-V1's timing optimization represented locally optimal decisions that were globally suboptimal, a failure mode characteristic of LLMs operating without sufficient domain context. This suggests that prompt engineering for cyber-physical attacks must go beyond protocol-level instructions to include system dynamics.

Second, domain model awareness is essential: AI-V2 succeeded only after receiving explicit knowledge of power accumulation dynamics, demonstrating that cyber-physical LLM attackers require grounding in physical models, not just protocol knowledge. For defenders, this implies that attacks requiring multi-step physical reasoning may be detectable through their characteristic patterns of information gathering.

Third, the framework enables systematic attacker development: the V1-to-V2 evolution illustrates how LLM-GridEval supports iterative refinement of attacker policies within a controlled environment, with each variant using identical primitives and constraints to isolate the impact of policy changes. This capability is essential for rigorous security evaluation where confounding variables must be minimized.

The current experiments represent initial validation under favorable attack conditions. Future work will evaluate longer campaigns (1--2 hours), harder grid configurations, and additional attack primitives to more fully characterize the evaluation validity gap.


\section{Limitations and Future Work}

The current evaluation has several limitations that scope the conclusions. First, the proof-of-concept experiments use a single operating point (simulation hour 7) with favorable attack conditions; scenarios with lower base load, higher thresholds, or faster controller response would better differentiate attacker capabilities and more clearly reveal the evaluation validity gap. Second, the 5-minute experiment window captures only short-term attack-response dynamics; longer campaigns spanning 1--2 hours would test resource management strategies and reveal whether the efficiency advantages of adaptive attackers compound over time. Third, the threshold-based defender is simpler than production systems, which may incorporate anomaly detection, predictive control, or coordinated protection schemes. Fourth, the current primitive catalog is limited to EV capacity manipulation; real adversaries may combine load-altering attacks with false data injection or communication disruption. Finally, the software co-simulation cannot fully capture the timing constraints and communication latencies of real power system hardware.

Several extensions are planned: (1) replacing scripted policies with fully autonomous LLM agents using multi-agent architectures similar to PentestGPT and AttackLLM \cite{deng2024pentestgpt,attackllm}, with retrieval-augmented generation supplying power system models and attack templates; (2) integrating data-driven and physics-informed intrusion detection systems to quantify $\Delta E(M)$ for detection mechanisms and enable controlled comparisons between RL-based and LLM-based attackers \cite{aoufi2020fdiasurvey,sahani2023smartgridids}; (3) expanding the primitive catalog to include false data injection, malicious SCADA commands, and communication disruption for multi-stage attack evaluation; (4) using LLM-GridEval to generate diverse synthetic attack traces that capture adaptive adversary behaviors for IDS benchmarking, complementing established static datasets; and (5) validating the framework on hardware-in-the-loop (HIL) testbeds such as OPAL-RT to evaluate LLM-based attacks against real-time power system simulators with realistic timing constraints and communication latencies. We also plan to release the framework as open-source software to support reproducible security research in the smart grid community.

\section{Conclusion}

This paper introduced LLM-GridEval, a framework that embeds adaptive, tool-using LLM attackers into HELICS-based power grid co-simulations through a schema-constrained primitive interface. Our proof-of-concept evaluation demonstrated that timing intelligence alone is insufficient for effective attacks; LLM-based adversaries require explicit domain model awareness to outperform random baselines. The V1 to V2 attacker evolution illustrates how the framework enables systematic study of adaptive attack strategies under controlled conditions. These results underscore that static evaluations may underestimate adversarial risk in API-driven, EV/DER-rich grids, motivating further research on LLM-based attackers, advanced defenders, and richer attack scenarios within physically grounded evaluation environments.



% \section*{Acknowledgment}

% \section*{References}

\bibliographystyle{IEEEtran}
\bibliography{ref}


%\bibliographystyle{IEEEtranS}
%\bibliography{ref}



\end{document}
